% Generated draft: requires mathematical and source review.
\documentclass[11pt]{article}
\usepackage[margin=25mm]{geometry}
\usepackage{fontspec}
\setmainfont{Latin Modern Roman}
\usepackage{amsmath,amssymb,mathtools,unicode-math}
\setmathfont{Latin Modern Math}
\usepackage{enumitem,needspace,xurl,hyperref}
\setlength{\parindent}{0pt}
\setlength{\parskip}{5pt}
\setlength{\emergencystretch}{3em}
\newcommand{\N}{\mathbb{N}}
\newcommand{\Z}{\mathbb{Z}}
\newcommand{\Q}{\mathbb{Q}}
\newcommand{\R}{\mathbb{R}}
\newcommand{\C}{\mathbb{C}}
\newcommand{\F}{\mathbb{F}}
\newcommand{\A}{\mathbb{A}}
\newcommand{\PP}{\mathbb{P}}
\newcommand{\E}{\mathbb{E}}
\newcommand{\eps}{\varepsilon}
\newcommand{\dd}{\,\mathrm{d}}
\begin{document}
\section*{UnsolvedMath ID: 20003086}
\textbf{Source problem code:} catalogue entry
\subsection{Definitions and mathematical statement}
This record is contextual prose, not a separately stated mathematical problem. Its supplied text says that Problem 1.1 is motivated by Royden's proof that, roughly speaking, the isometry group of Teichm\"uller space with its Teichm\"uller metric is the mapping class group. The aid reports that this sentence immediately follows the genuine Problem 1.1 in the AIM PDF and that the long dot run is a printed separator, not omitted mathematics. The genuine problem is not reproduced in this record. Thus no independent conjecture or exact unrestricted version of Royden's theorem can be extracted from this sentence alone; surface-complexity and effective-group qualifications are not given.

The aid appends a formal reduction for a different, Thurston-metric rigidity program. In a finite-type surface setting where the definitions apply, the infinitesimal asymmetric Thurston norm is
\[
 \|v\|_X=\sup_{\alpha}
       \frac{d\ell_\alpha|_X(v)}{\ell_\alpha(X)},
\]
with $\alpha$ ranging over essential simple closed curves. Mapping classes act equivariantly on curve lengths, so their derivatives identify the normed tangent spaces at $X$ and $gX$. Consequently any recovery from the isometry class of this normed space can distinguish at most the effective extended mapping-class orbit, not a unique marked point.

Here is the conditional globalization reduction reported in the aid. Let $G$ be the effective extended mapping class group, let $\mathcal T_{\mathrm{free}}$ be the locus with trivial $G$-stabilizer, and use its properly discontinuous free action, so
\[
 q:\mathcal T_{\mathrm{free}}\longrightarrow
       \mathcal T_{\mathrm{free}}/G
\]
is a covering map. Let $U\subset\mathcal T_{\mathrm{free}}$ be connected and let $F:U\to\mathcal T_{\mathrm{free}}$ be continuous. If $F(X)\in G\cdot X$ for every $X\in U$, equivalently $q\circ F=q|_U$, then there is a single $g\in G$ with $F=g|_U$. Freeness makes the pointwise group element unique and the covering property makes it locally constant; connectedness gives one element on all of $U$.

In particular, suppose infinitesimal norm rigidity is available in the chosen surface class, in the form that a linear Thurston-norm isometry between tangent spaces forces their basepoints into the same $G$-orbit. A $C^1$ local Thurston-metric isometric embedding on such a connected free domain then has this pointwise orbit property through its derivative, and the covering reduction makes it one mapping class on the domain. This is a conditional synthesis, not the original contextual sentence. Differentiability, stabilizer loci, low-complexity exceptions and extension beyond the free locus need separate input.

The aid states that the surrounding infinitesimal and local rigidity program has been solved by Pan and by Huang--Ohshika--Papadopoulos, without supplying precise theorem statements, surface hypotheses, or full references. That reported literature status does not turn this source fragment into an independent open problem.
\subsection{Short English statement}
This entry is the motivational sentence following another AIM problem, so it supplies no independent question. The attached aid discusses a conditional rigidity reduction: equivariant norm data determine a mapping-class orbit, and a covering argument turns pointwise orbit agreement into one mapping class on a connected domain without stabilizers.
\subsection{Sources}
\begin{itemize}
\item Local UnsolvedMath registry, permanent ID $20003086$, source code catalogue entry; canonical location \texttt{aim-topology-notes.json}, \texttt{notes[173]}.
\item American Institute of Mathematics, \emph{Lipschitz metric on Teichmueller space}, motivational prose following Problem 1.1, \url{https://aimath.org/pastworkshops/lipschitzteichproblems.pdf}. The supplied aid reports the contextual placement and printed separator; no PDF was accessed for this draft. The sentence attributes its motivation to Royden.
\item Attached research attempt 1 and summary name Pan and Huang--Ohshika--Papadopoulos without full bibliographic citations. They mark the attempt context-only and the source not a problem, despite the registry's open label. A proved/low-novelty label is accompanied by the explicit machine-generated, non-peer-reviewed research-aid caveat requiring independent verification before citation.
\end{itemize}
\end{document}
